\documentclass[11pt]{scrartcl}

\usepackage[stixtwo]{fontsetup}
\usepackage{amsmath}
\usepackage{xfrac}

\begin{document}
\begin{enumerate}
	\item \textbf{Szöveges megoldás}

	      A számlálóban a leggyorsabban növekvő tag a $3 \cdot 3^{2n - 1}
		      = 9^n$, hiszen
	      \[
		      n \ln n \ll n^2 \ll 9^n
		      \text.
	      \]
	      A nevezőben a leggyorsabban növekvő tag a $3^{n + 2} = 9 \cdot 3^n$,
	      hiszen
	      \[
		      n^4 \ll 3^n
		      \text.
	      \]
	      Mivel a számlálóban a $9^n$, a nevezőben pedig a $c \cdot 3^n$ a
	      domináns tag, és $9^n$ sokkal gyorsabban nő, mint $3^n$, ezért a
	      határérték végtelen lesz:
	      \[
		      \lim_{n \to \infty} \frac{
			      n \, \ln(n) + 3 \cdot 3^{2n - 1}
		      }{
			      3^{n + 2} + 6 n^4
		      } = \infty
		      \text.
	      \]

	      \textbf{Másik megoldás}

	      $3 \cdot 3^{2n - 1} = 9^n$ és $3^{n + 2} = 9 \cdot 3^n$, így
	      \[
		      \frac{n \ln(n) + 9^n}{9 \cdot 3^n + 6n^4}
		      = \frac{\frac{n \ln(n)}{3^n} + 3^n}{9 + \frac{6n^4}{3^n}}
		      \to \infty
		      \text.
	      \]

	\item Behelyettesítéssel $(\infty - \infty)$ típusú határozatlan alakot
	      kapunk. Alkalmazzuk a gyökös kifejezések racionalizálását, vagyis
	      bővítsük a kifejezést a konjugáltjával:
	      \[
		      \sqrt{n^2 + 3n + 2} - n
		      = \frac{
			      (n^2 + 3n + 2) - n^2
		      }{
			      \sqrt{n^2 + 3n + 2} + n
		      }
		      = \frac{3n + 2}{\sqrt{n^2 + 3n + 2} + n}
		      \text.
	      \]
	      A határérték tehát:
	      \[
		      = \frac{3n + 2}{\sqrt{n^2 + 3n + 2} + n}
		      = \frac{3 + \sfrac{2}{n}}{\sqrt{1 + \sfrac{3}{n} + \sfrac{2}{n^2}} + 1}
		      \to \frac{3}{1 + 1} = \frac{3}{2}
		      \text.
	      \]

	\item \textbf{Számláló és nevező együtt kezelésével}
	      \[
		      \left( \frac{3n + 2}{7 + 3n} \right)^{n + 2}
		      = \left( \frac{3n + 7 - 5}{7 + 3n} \right)^{n + 2}
		      = \left[
			      \left( 1 + \frac{-5}{7 + 3n} \right)^{7 + 3n}
			      \right]^{\frac{n + 2}{7 + 3n}}
		      \to \left[e^{-5}\right]^{\frac{1}{3}} = e^{-\sfrac{5}{3}}
	      \]

	      \textbf{Számláló és nevező külön kezelésével}
	      \[
		      \left( \frac{3n + 2}{7 + 3n} \right)^{n + 2}
		      = \frac{\left(1 + \frac{2/3}{n}\right)^n}{\left(1 + \frac{7/3}{n}\right)^n}
		      \cdot \left( \frac{3n + 2}{3n + 7} \right)^2
		      \to \frac{e^{\sfrac{2}{3}}}{e^{\sfrac{7}{3}}} \cdot 1
		      = e^{-\sfrac{5}{3}}
	      \]
\end{enumerate}
\end{document}
